\documentclass[12pt,a4paper]{article}

\usepackage[utf8]{inputenc}
\usepackage[T1]{fontenc}
\usepackage[brazil]{babel}
\usepackage{geometry}
\usepackage{hyperref}
\usepackage{enumitem}
\usepackage{titlesec}
\usepackage{graphicx}
\usepackage{indentfirst}

\geometry{
  top=2.5cm,
  bottom=2.5cm,
  left=3cm,
  right=2.5cm
}

\hypersetup{
  colorlinks=true,
  linkcolor=blue,
  urlcolor=blue
}

\begin{document}

%--------------------------------------------------------------
% CAPA
%--------------------------------------------------------------
\begin{titlepage}
  \centering
  {\large UNIVERSIDADE ESTADUAL DO NORTE FLUMINENSE}\\[0.3cm]
  {\large LABORATÓRIO DE ENGENHARIA E EXPLORAÇÃO DE PETRÓLEO}\\[0.3cm]
  {\large CENTRO DE CIÊNCIA E TECNOLOGIA}\\[1.5cm]
  {\Large \textbf{MICRO PROJETO DE ENGENHARIA}}\\[0.5cm]
  {\Large \textbf{DESENVOLVIMENTO DO SOFTWARE}}\\[0.5cm]
  {\Large \textbf{XXXX}}\\[1.5cm]
  {\normalsize DISCIPLINA: Introdução ao Projeto de Engenharia}\\[0.3cm]
  {\normalsize Setor de Modelagem Matemática Computacional}\\[1.5cm]
  \begin{flushleft}
    \textbf{Equipe Clientes:}\\
    nomes\\[0.8cm]
    \textbf{Equipe Desenvolvedores:}\\
    nomes\\[0.8cm]
    \textbf{Revisor Técnico:}\\
    Prof. André Duarte Bueno
  \end{flushleft}
  \vfill
  {\normalsize MACAÉ -- RJ}\\[0.2cm]
  {\normalsize Janeiro -- 2026}
\end{titlepage}

%--------------------------------------------------------------
% SEÇÃO: DIDÁTICA
%--------------------------------------------------------------
\section*{Didática}

O objetivo didático é implementar a Aprendizagem Baseada em Projetos (PBL) e simulação de papéis (\textit{Role-Playing}).

Simular uma situação real do relacionamento entre clientes e empresas de engenharia.

A dinâmica emula o ciclo inicial de engenharia de requisitos e a comunicação (frequentemente ruidosa) entre:
\begin{itemize}
  \item ``O Cliente'' que entende da física/fenômeno do problema, normalmente um geólogo, geofísico, engenheiro.
  \item ``O Engenheiro de Software/Projetista'', que fará a modelagem do sistema de engenharia.
\end{itemize}

\subsection*{Limites\footnote{O limite tem fins didáticos e é estabelecido apenas para que o problema seja resolvido em curto espaço de tempo. Problemas reais são bem mais complexos e detalhados.} Estritos Recomendados (O Escopo)}

Para garantir que as duplas consigam entregar os diagramas no prazo, os alunos clientes devem preencher um ``Template de Especificação'' contendo, no máximo, os seguintes itens:

\begin{itemize}
  \item \textbf{Tema/Problema de Engenharia:}
  \begin{itemize}
    \item Restrito a um cálculo ou simulação simplificada (ex: coleta e análise de amostras catalogadas; processamento de dados de um ensaio de campo; cálculo básico de propriedades de rochas; análise de dados de perfis; análise geofísica; processamento de dados de fluidos; análise de dados de reservatório; simulação simples de sistemas simples aplicados a engenharia de reservatório, poço, elevação e escoamento).
  \end{itemize}

  \item \textbf{Atores (Usuários):}
  \begin{itemize}
    \item Máximo de 2
    \item (ex: Usuário, Operador de Laboratório, Sistema de Aquisição de Dados, Sistema Gerador de Gráficos).
  \end{itemize}

  \item \textbf{Requisitos Funcionais (RF):}
  \begin{itemize}
    \item 3 a 4 requisitos.
    \item Ex: RF1 -- O sistema deve ler um arquivo com as leituras de pressão e vazão. RF2 -- O sistema deve calcular a permeabilidade. RF3 -- O sistema deve gerar um relatório binarizado de falhas.
  \end{itemize}

  \item \textbf{Requisitos Não Funcionais (RNF):}
  \begin{itemize}
    \item 1 a 2 requisitos, focados em engenharia de software.
    \item Ex: RNF1 -- A modelagem deve prever o encapsulamento estrito visando codificação em C++ moderno. RNF2 -- O tempo de processamento das matrizes não deve exceder 2 segundos.
  \end{itemize}

  \item \textbf{Casos de Uso Principais:}
  \begin{itemize}
    \item Exigir apenas 1 ou 2 cenários (O ``Caminho Feliz'' do cálculo e 1 fluxo de exceção, como ``Dados de entrada inconsistentes'').
  \end{itemize}
\end{itemize}

%--------------------------------------------------------------
% SEÇÃO: METODOLOGIA
%--------------------------------------------------------------
\section*{Metodologia}

A Figura~\ref{fig:metodologia} apresenta a metodologia a ser utilizada no desenvolvimento do sistema (modelos).

\begin{itemize}
  \item Para facilitar disponibilizamos um repositório modelo deste micro-projeto de engenharia. O mesmo pode ser clonado:
  \begin{itemize}
    \item \url{https://github.com/ldsc/LDSC-ProjetoEngenharia-3-MicroProjetoSoftware-TituloProjeto}
  \end{itemize}
\end{itemize}

\begin{figure}[h!]
  \centering
  % Substituir pelo caminho da imagem exportada do PDF, se disponível:
  % \includegraphics[width=\textwidth]{metodologia.png}
  \fbox{\parbox{0.9\textwidth}{\centering [Diagrama da Metodologia -- inserir imagem]}}
  \caption{Metodologia utilizada no desenvolvimento do sistema}
  \label{fig:metodologia}
\end{figure}

%--------------------------------------------------------------
% CONTRATO
%--------------------------------------------------------------
\newpage
\begin{center}
  {\Large \textbf{Contrato de Prestação de Serviços de\\[0.3cm] Engenharia de Software}}
\end{center}

\vspace{0.5cm}

Pelo presente instrumento particular, de um lado, doravante denominado simplesmente CONTRATANTE (Grupo Cliente), e, de outro lado, doravante denominado simplesmente CONTRATADA (Grupo Executor), celebram o presente Contrato de Prestação de Serviços de Engenharia de Software, mediante as cláusulas e condições a seguir estipuladas.

\begin{itemize}[leftmargin=*]

  %--- CLÁUSULA 1 ---
  \item \textbf{CLÁUSULA 1 -- DO OBJETO}
  \begin{itemize}
    \item O presente contrato tem por objeto a especificação, modelagem UML e futuro desenvolvimento de um software aplicativo voltado à resolução de um problema específico de Engenharia de Petróleo, em conformidade estrita com as metodologias apresentadas na Figura~\ref{fig:metodologia} e nas referências \cite{Rumbaugh1994, Bueno2003, Blaha2006}.
  \end{itemize}

  %--- CLÁUSULA 2 ---
  \item \textbf{CLÁUSULA 2 -- ESCOPO DO PROBLEMA DE ENGENHARIA}
  \begin{itemize}
    \item \textbf{Objetivo:}
    \begin{itemize}
      \item Ex: O PercolaEdu 2D é um software educativo interativo desenvolvido para auxiliar estudantes do ensino médio e superior (Física, Matemática, Ciência da Computação, Engenharia) a compreenderem os conceitos de probabilidade, transição de fase, conectividade e teoria dos grafos através do modelo de percolação de sítios em malhas bidimensionais (2D).
    \end{itemize}

    \item \textbf{Público-Alvo:}
    \begin{itemize}
      \item Alunos do ensino médio avançado.
      \item Estudantes de graduação em Ciências Exatas.
      \item Professores que necessitam de ferramentas visuais para demonstrar fenômenos críticos e estatísticos.
    \end{itemize}

    \item \textbf{Especificação:}
    \begin{itemize}
      \item Ex: O software educativo PercolaEdu 2D, com licença GPL -- livre, deve ser multiplataforma, usar linguagem de alto desempenho (C++) e ter uma interface em modo terminal -- CLI. O simulador deve permitir a entrada de dados da rede de percolação, suas dimensões $M_x$, $N_y$, e $p$ a probabilidade de ativação de cada sítio. Podendo, opcionalmente, informar o modelo de distribuição a ser utilizado para determinar a probabilidade $p$ (distribuição uniforme, normal, Bernoulli, T-Student) e o modelo de vizinhança pode ser de faces diretas ou faces e cantos. Deve ser possível gerar redes únicas para estudo de um caso específico (um único valor de $p$) ou situações em que o valor da probabilidade $p$ cresce até que ocorra a percolação. Os resultados devem incluir a informação se ocorre ou não a percolação, ou seja, se existe um caminho contínuo de sítios ocupados conectando o topo à base (ou a esquerda à direita) da malha. As imagens resultantes devem ser salvas e opcionalmente visualizadas. Os algoritmos de agrupamento devem incluir a busca profunda e Hoshen-Kopelman.
    \end{itemize}

    \item \textbf{Cenários:}
    \begin{itemize}
      \item A fim de delimitar a fronteira de interação entre os atores e o sistema, ficam estabelecidos os cenários principais de uso abaixo descritos:

      \item \textbf{Cenário 1:}
      \begin{itemize}
        \item Ex: O software permitirá que o usuário crie uma malha (\textit{grid}) de tamanho $M_x$, $N_y$. Cada ``sítio'' (célula) da malha será ocupado com uma probabilidade $p$ ou vazio com probabilidade $1-p$. O sistema deve identificar visualmente os ``clusters'' (agrupamentos de sítios vizinhos ocupados) e determinar se houve a percolação, ou seja, se existe um caminho contínuo de sítios ocupados conectando o topo à base (ou a esquerda à direita) da malha.
      \end{itemize}

      \item \textbf{Cenário 2:}
      \begin{itemize}
        \item Ex: O sistema não percolou e é necessário aumentar a probabilidade $p$. Ou seja, é iterativo. Mostra o passo a passo.
      \end{itemize}
    \end{itemize}

    \item \textbf{Requisitos:}

    \medskip
    \textbf{RF -- Requisito Funcional}

    \begin{itemize}
      \item \textbf{RF01}
      \begin{itemize}
        \item \textit{Nome:} Ex: Configuração da Malha
        \item \textit{Descrição:} O usuário deve poder definir o tamanho da malha quadrada inserindo um valor $N$ (ex: de 10 a 500).
        \item \textit{Prioridade:} Alta
      \end{itemize}

      \item \textbf{RF02}
      \begin{itemize}
        \item \textit{Nome:} Ex: Definição da Probabilidade ($p$)
        \item \textit{Descrição:} O usuário deve poder ajustar a probabilidade de ocupação dos sítios usando um controle deslizante (\textit{slider}) ou caixa de texto, com valores entre 0,000 e 1,000.
        \item \textit{Prioridade:} Alta
      \end{itemize}

      \item \textbf{RF03}
      \begin{itemize}
        \item \textit{Nome:} Ex: Geração Aleatória
        \item \textit{Descrição:} O sistema deve gerar a rede preenchendo os sítios aleatoriamente com base na probabilidade $p$ definida no RF02.
        \item \textit{Prioridade:} Alta
      \end{itemize}

      \item \textbf{RF04}
      \begin{itemize}
        \item \textit{Nome:} Ex: Identificação de Clusters
        \item \textit{Descrição:} O algoritmo (ex: Union-Find ou Busca em Profundidade -- DFS) deve agrupar sítios adjacentes ortogonalmente. Cada cluster distinto deve receber uma cor diferente na interface gráfica.
        \item \textit{Prioridade:} Alta
      \end{itemize}

      \item \textbf{RF05}
      \begin{itemize}
        \item \textit{Nome:} Ex: Detecção de Percolação
        \item \textit{Descrição:} O sistema deve verificar automaticamente se o maior cluster conecta as bordas opostas (topo-base ou esquerda-direita) e destacar esse ``cluster percolante'' com uma cor especial (ex: vermelho brilhante).
        \item \textit{Prioridade:} Alta
      \end{itemize}
    \end{itemize}

    \medskip
    \textbf{RNF -- Requisito Não Funcional}

    \begin{itemize}
      \item \textbf{RNF01 -- Plataforma:}
      \begin{itemize}
        \item O software deve ser uma aplicação modo terminal multiplataforma.
      \end{itemize}

      \item \textbf{RNF02 -- Tecnologias Sugeridas:}
      \begin{itemize}
        \item Ex: Gerar as imagens em formato portável entre plataformas.
      \end{itemize}

      \item \textbf{RNF03 -- Desempenho:}
      \begin{itemize}
        \item Ex: A geração e a verificação de percolação de uma malha de $500 \times 500$ devem ocorrer em menos de 1 segundo.
      \end{itemize}

      \item \textbf{RNF04 -- Usabilidade:}
      \begin{itemize}
        \item Ex: A interface deve seguir princípios de design limpo (UI/UX), utilizando textos claros, disponibilizando manuais simplificados.
      \end{itemize}
    \end{itemize}

    \medskip
    \textbf{RN -- Regras de Negócio}

    \begin{itemize}
      \item \textbf{RN01 -- Vizinhança:}
      \begin{itemize}
        \item Ex: Para a formação de um cluster, a conexão entre sítios deve considerar apenas a vizinhança de von Neumann (vizinhos diretos: cima, baixo, esquerda, direita). Diagonais não formam conexões contínuas.
      \end{itemize}

      \item \textbf{RN02 -- Limites de Entrada:}
      \begin{itemize}
        \item Ex: O sistema não deve aceitar valores de $M_x$, $N_y$ menores que 2 ou maiores que um limite seguro estabelecido pelo hardware do cliente (ex: 10000).
      \end{itemize}
    \end{itemize}
  \end{itemize}

  %--- CLÁUSULA 3 ---
  \item \textbf{CLÁUSULA 3 -- PRAZOS E CRONOGRAMA (SPRINT)}

  As partes acordam o seguinte cronograma de entregas técnicas, sujeito à validação da CONTRATANTE:

  \begin{itemize}
    \item \textbf{Dia 1: Concepção Preliminar e Proposta (Em Sala e Em Casa)}
    \begin{itemize}
      \item \textit{Cliente:}
      \begin{itemize}
        \item Elabora e entrega o documento ``ESCOPO DO PROBLEMA DE ENGENHARIA''\footnote{Problema vinculado ao curso de engenharia}, detalhando Objetivo, Público-Alvo, Especificação, Cenários, Requisitos.
      \end{itemize}
      \item \textit{Executor (Pré-venda):}
      \begin{itemize}
        \item Monta a Proposta Técnica contendo a experiência da equipe, uma elaboração preliminar do problema (o que entenderam), primeira versão do Diagrama de Casos de Uso (Concepção), esboço do Diagrama de Classes (Domínio) e primeiro rascunho do Diagrama de Sequência. Orçamento e Cronograma Preliminar.
        \item Entrega: Via e-mail ou repositório temporário para avaliação do cliente.
      \end{itemize}
      \item \textit{Cliente (Em casa):}
      \begin{itemize}
        \item Revisão da proposta enviada pelo executor e preparação de dúvidas.
        \item Verificar se houve um claro entendimento do problema a ser resolvido.
      \end{itemize}
    \end{itemize}

    \item \textbf{Dia 2: Elaboração, Contratação e Análise OO (Reunião Formal e Em Casa)}
    \begin{itemize}
      \item \textit{Reunião Formal:}
      \begin{itemize}
        \item Executor apresenta competências da equipe (histórico de projetos desenvolvidos -- marketing); soluções arquiteturais desenvolvidas (modelos UML), orçamento e cronograma preliminar.
        \item Cliente analisa a proposta, verifica e aprofunda os requisitos técnicos (cálculos, física, química, engenharia) e valida a viabilidade técnica.
        \item Negociação e Assinatura deste Contrato\footnote{Simulem de fato; para assinar usem o gov.br. Este contrato não tem valor legal, é um exemplo de sala de aula.}.
      \end{itemize}
      \item \textit{Executor (Em casa):}
      \begin{itemize}
        \item Configuração oficial do Repositório Git/GitHub (aplicação de versionamento inicial). Clonar o repositório modelo, adicionar membros (clientes e equipe de desenvolvimento).
        \item Execução da fase de Elaboração: Estudo do problema. Montagem do Diagrama de Pacotes / Assuntos.
        \item Consolidação da Análise OO: Montagem dos Diagramas de Casos de Uso Geral e Específicos, Diagrama de Assuntos e Pacotes, Diagrama de Classes, Diagramas de Sequência (cenários principais), Diagrama de Comunicação e Diagrama de Máquina de Estado.
        \item Documentação para apresentação formal.
        \item Entrega: Registro das atividades (checklist, commit) e envio dos dados via repositório do projeto.
      \end{itemize}
    \end{itemize}

    \item \textbf{Dia 3: Projeto do Sistema e Revisão (Reunião Formal e Em Casa)}
    \begin{itemize}
      \item \textit{Reunião Formal:}
      \begin{itemize}
        \item Executor apresenta o modelo de Análise OO consolidado.
        \item Cliente comenta, critica e solicita correções formais.
      \end{itemize}
      \item \textit{Executor (Em casa):}
      \begin{itemize}
        \item Definições de Projeto do Sistema: Definição formal de Arquitetura, Bibliotecas (Polinômios/Estatística) e encapsulamento em C++\footnote{O cliente exige C++ porque quer alto desempenho e economia de energia (é um ambientalista).}.
        \item Projeto OO: Criação do Diagrama de Componentes e Diagrama de Implantação (\textit{hardware}/\textit{cluster}).
        \item Revisão e refatoração de todos os diagramas UML com base na arquitetura definida.
      \end{itemize}
    \end{itemize}

    \item \textbf{Dia 4: Entrega da Modelagem Abrangente e Montagem do Backlog -- Lista de Características (Reunião Formal)}
    \begin{itemize}
      \item \textit{Reunião Formal:}
      \begin{itemize}
        \item Executor apresenta a documentação final do Ciclo de Concepção e Análise. Inclui documento formal com descrição de todas as etapas desenvolvidas. O Executor extrai dos modelos UML aprovados e entrega a Lista de Características (\textit{Backlog} do Produto), fatiando o escopo em entregas $c_1, c_2, \ldots, c_n$, habilitando o início imediato do Ciclo de Planejamento e Construção Iterativa.
        \item Cliente: comenta, critica, realiza o aceite técnico formal da documentação entregue e decide sobre a continuidade da mesma empresa/equipe para os ciclos de planejamento e construção\footnote{Se o resultado é satisfatório, contrate para a próxima etapa; se for mediano ou ruim, seja gentil, mas não dê continuidade -- procure outra equipe. Nada de tapinha nas costas; se ficou insatisfeito, procure outra equipe. Haja com profissionalismo, afinal de contas você tem um chefe.}.
        \item Marco Final (\textit{Deliverable}): Documento final com todos os diagramas e o Backlog -- Lista de Características. Tudo atualizado e armazenado no repositório do projeto.
      \end{itemize}
    \end{itemize}
  \end{itemize}

  %--- CLÁUSULA 4 ---
  \item \textbf{CLÁUSULA 4 -- CUSTOS E CONDIÇÕES DE PAGAMENTO}
  \begin{itemize}
    \item A remuneração pelo projeto baseia-se no esforço de engenharia empregado na especificação, modelagem, testes lógicos e documentação, eventualmente alguma codificação.
    \item Valor da Hora de Engenharia de Software: R\$ 350,00.
    \item Custo Total Estimado para o MVP (\textit{Minimum Viable Product}): [Definir valor baseado no esforço de modelagem -- sugerido 40h de engenharia].
    \item Os recursos computacionais físicos para compilação e teste final são de inteira responsabilidade da CONTRATADA.
  \end{itemize}

  %--- CLÁUSULA 5 ---
  \item \textbf{CLÁUSULA 5 -- REQUISITOS TECNOLÓGICOS E EFICIÊNCIA ENERGÉTICA}
  \begin{itemize}
    \item Fica estabelecido que a implementação final da arquitetura modelada deverá ser conduzida na linguagem C++ (padrões modernos C++26 ou mais recente). Esta exigência contratual é inegociável e visa garantir não apenas a máxima performance de execução (cálculos massivos de engenharia), mas também a eficiência energética e ecológica, mitigando o desperdício de processamento, consumo de energia e recursos naturais frequentemente associados a linguagens interpretadas e antiecológicas.
  \end{itemize}

  %--- CLÁUSULA 6 ---
  \item \textbf{CLÁUSULA 6 -- SIGILO E PROPRIEDADE INTELECTUAL}
  \begin{itemize}
    \item A CONTRATADA compromete-se a manter absoluto sigilo sobre os dados de campo, dados laboratoriais e amostras de laboratório fornecidos pela CONTRATANTE que tenham a marcação SIGILOSA.
    \item A modelagem UML (diagramas e arquivos LyX) e o código-fonte desenvolvido em C++ serão disponibilizados sob a licença GPL no site \url{http://www.github.com/ldsc} em repositório específico a ser criado.
  \end{itemize}

  %--- CLÁUSULA 7 ---
  \item \textbf{CLÁUSULA 7 -- FORO}
  \begin{itemize}
    \item Fica eleito o foro da sala de aula, Juiz Prof. André D. Bueno, para dirimir quaisquer controvérsias oriundas da execução deste contrato de engenharia.
  \end{itemize}

\end{itemize}

\vspace{1cm}
Macaé -- RJ, \underline{\hspace{1cm}} de \underline{\hspace{3cm}} de 202\underline{\hspace{0.5cm}}.

\vspace{2cm}
\noindent
\begin{tabular}{p{0.45\textwidth}p{0.45\textwidth}}
\hline
\multicolumn{2}{l}{\textbf{CONTRATANTE (Engenheiro Cliente)}}\\[0.2cm]
\multicolumn{2}{l}{Grupo:}\\[1cm]
Engenheiro \underline{\hspace{3cm}} & Assinatura: \underline{\hspace{3cm}} \\[1cm]
Engenheiro \underline{\hspace{3cm}} & Assinatura: \underline{\hspace{3cm}} \\
\end{tabular}

\vspace{1.5cm}
\noindent
\begin{tabular}{p{0.45\textwidth}p{0.45\textwidth}}
\hline
\multicolumn{2}{l}{\textbf{CONTRATADA (Engenheiro Executor)}}\\[0.2cm]
\multicolumn{2}{l}{Grupo:}\\[1cm]
Engenheiro \underline{\hspace{3cm}} & Assinatura: \underline{\hspace{3cm}} \\[1cm]
Engenheiro \underline{\hspace{3cm}} & Assinatura: \underline{\hspace{3cm}} \\
\end{tabular}

\vspace{1.5cm}
\noindent\rule{0.6\textwidth}{0.4pt}\\
Testemunha / Gestor de Projeto (Professor André Bueno)

%--------------------------------------------------------------
% REFERÊNCIAS
%--------------------------------------------------------------
\newpage
\begin{thebibliography}{9}

\bibitem{Blaha2006}
BLAHA, M.; RUMBAUGH, J.
\textit{Modelagem e Projetos Baseados em Objetos com UML~2}.
Rio de Janeiro: Campus, 2006.

\bibitem{Bueno2003}
BUENO, A. D.
\textit{Programação Orientada a Objeto com C++ -- Aprenda a Programar em Ambiente Multiplataforma com Software Livre}.
São Paulo: Novatec, 1.~ed., 2003.

\bibitem{Rumbaugh1994}
RUMBAUGH, J.; BLAHA, M.; PREMERLANI, W.; EDDY, F.; LORENSEN, W.
\textit{Modelagem e Projetos Baseados em Objetos}.
Rio de Janeiro: Campus, 1994.

\end{thebibliography}

\end{document}
